\section{Introduction}

The excess--charge problem asks how many electrons an atomic nucleus of charge $Z$ can bind.  In the nonrelativistic Schr\"odinger model one expects only a bounded number of excess electrons, independently of $Z$.  A classical theorem of Lieb gives $N_c(Z)<2Z+1$ \cite{Lieb1984}; Nam later improved the coefficient to $1.22$ with an explicit $O(Z^{1/3})$ correction \cite{Nam2012}.  The conjectural bound remains $N_c(Z)\le Z+C$ for a universal constant $C$; see, for example, \cite{Nam2013Review}.

A recent refinement by Hundertmark, Pattakos and Schulz (HPS) develops the weighted multiplier method substantially further and proves
\[
 N_c(Z)<1.1185\,Z+O(Z^{1/3}),
\]
with an explicit lower--order term \cite{HPS2025}.  A central object in their argument is the mean--field constant
\begin{equation}\label{eq:def-beta-s}
 \beta_s:=\inf_{\mu}
 \frac{\displaystyle
 \iint_{\R^3\times\R^3}
 \frac{|x|^s+|y|^s}{2|x-y|}\,\dd\mu(x)\dd\mu(y)}
 {\displaystyle\int_{\R^3}|x|^{s-1}\,\dd\mu(x)},
\end{equation}
where the infimum is taken over the admissible probability measures used in \cite{HPS2025}.  HPS prove that for $1\le s\le3$ the infimum can be restricted to radial measures.  For $s=3$ they then use a pointwise estimate on the radial kernel.  This yields the explicit lower bound
\[
 \beta_3\ge 0.8941074569\ldots,
\]
whose reciprocal is the source of the leading coefficient $1.1185$ in their numerical theorem.

The first purpose of the present paper is to solve the radial variational problem at $s=3$ exactly.

\begin{theorem}[Exact value of $\beta_3$]\label{thm:main}
Let $\beta_3$ be the HPS mean--field constant defined by \eqref{eq:def-beta-s}.  Set
\[
 a_*:=\frac{1}{\sqrt2}\arctan\sqrt2.
\]
Then
\begin{equation}\label{eq:beta-exact}
 \boxed{\displaystyle
 \beta_3=\frac{3\ee^{2a_*}}{1+3\ee^{2a_*}}}
 =0.9205346822904167\ldots .
\end{equation}
In particular,
\[
 \beta_3^{-1}=1.086325175181735\ldots .
\]
A minimizing radial probability measure is supported on the shell
\[
 \ee^{-a_*}\le |x|\le \ee^{a_*}
\]
and has spatial density proportional to
\[
 |x|^{-4}\cos\!\bigl(\sqrt2\log|x|\bigr).
\]
Among radial admissible measures, the minimizer is unique up to the scaling invariance of \eqref{eq:def-beta-s}.  No uniqueness assertion is made here for minimizers in the full nonradial HPS class.
\end{theorem}

The second purpose is to propagate \Cref{thm:main} through the finite--$N$ part of the HPS argument.  Their cubic-weight estimates retain the actual mean--field constant $\beta_3$ symbolically; the pointwise lower bound enters only when the final coefficients are converted to explicit numbers.  Thus the present finite--$Z$ improvement is not obtained by replacing a decimal in the final HPS theorem: it comes from inserting the exact value \eqref{eq:beta-exact} into their symbolic chain and redoing the endpoint optimization.

\begin{theorem}[Sharpened finite--$Z$ excess--charge bound]\label{thm:finite-Z}
For every $Z\ge4$,
\begin{equation}\label{eq:finite-Z-main}
\boxed{\begin{aligned}
N_c(Z)<{}&\beta_3^{-1}Z+3.812\,Z^{1/3}+0.01294\\
&+0.18185\,Z^{-1/3}+0.01941\,Z^{-2/3},
\end{aligned}}
\end{equation}
where $\beta_3$ is the exact constant in \eqref{eq:beta-exact}.  Consequently,
\begin{equation}\label{eq:finite-Z-simple}
 \boxed{N_c(Z)<1.086326\,Z+3.90\,Z^{1/3}},\qquad Z\ge4.
\end{equation}
\end{theorem}

\subsection*{Significance of the result}

The improvement has a direct quantitative meaning.  The HPS leading coefficient $1.1185$ is produced by a lower estimate on the cubic mean--field constant, whereas the present computation replaces that estimate by the exact radial value.  The leading coefficient therefore drops to
\[
 1.086325175\ldots .
\]
Relative to the conjectural coefficient $1$, the distance from the HPS value is $0.1185$, while the exact radial computation removes
\[
 1.1185-1.086325175\ldots=0.032174824\ldots,
\]
or about $27\%$ of that leading-order gap.  This comparison is only about the linear coefficient; the $Z^{1/3}$ correction remains part of the finite--$Z$ theorem.  The point is that this gain is obtained without a new many--body inequality.  It isolates the amount of the HPS loss that came specifically from replacing a global mean--field optimization by a pointwise pairwise bound.

There is also a structural meaning.  The pointwise HPS estimate asks for the worst ratio associated with one pair of radii, but a probability measure cannot make every pair simultaneously realize that worst geometry.  Solving the full variational problem enforces this global compatibility.  The optimizer is not a singular pair configuration: it is an extended density supported on a finite annulus, with the explicit log--radial cosine profile in \Cref{thm:main}.  Thus the improvement is a genuinely collective effect inside the scalar radial mean--field problem.

Finally, the extremizer identifies the obstruction that remains after the scalar radial problem is solved exactly.  In physical radius the density is a compactly supported cosine modulation of $r^{-4}$, the critical power already suggested by the cubic weighted multiplier structure.  Consequently, any further reduction of the leading coefficient within a strengthened cubic-weight argument must use information not present in the scalar radial probability measure alone---for example fermionic occupation constraints or positive terms that are discarded in the purely mean--field reduction.  The exact profile therefore provides a concrete reference state for future stability estimates around the radial obstruction.

The proof of \Cref{thm:main} has three main steps.  First, Newton's theorem and the logarithmic change of variables $t=\log r$ transform the radial minimization into a one--dimensional maximization problem for the kernel
\[
 k(t)=\ee^{-|t|}-\ee^{-2|t|}.
\]
Second, the Euler--Lagrange obstacle condition determines a compactly supported cosine profile.  Third, a conditional negative--definiteness estimate on the extremal support rules out all competing radial maximizers.  The last step is where the exact solution becomes global within the radial problem rather than merely stationary.

The finite--$Z$ step deliberately keeps the HPS finite--configuration and weighted kinetic inequalities unchanged.  The new input is the exact solution of the mean--field variational problem, together with a finite--$Z$ reoptimization that tracks the admissible interval for $N/Z$.  One small endpoint issue must be handled explicitly: for $Z\ge4$, Lieb's bound gives
\[
 \frac NZ<2+\frac1Z\le\frac94,
\]
which is sharper than the coarser interval used in the original numerical optimization and is sufficient to locate the relevant supremum at the left endpoint $N/Z=\beta_3^{-1}$.

The resulting extremal profile also clarifies a numerical feature already visible in HPS: radial trial densities close to $|x|^{-4}$ are nearly optimal.  In logarithmic radius, the exact extremizer is a compactly supported cosine perturbation of precisely that critical power law.

The paper is organized as follows.  \Cref{sec:radial} records the radial reduction.  \Cref{sec:log} introduces the logarithmic variational problem.  \Cref{sec:extremizer} constructs the candidate maximizer and proves its obstacle condition.  \Cref{sec:global} proves global optimality for the radial problem.  \Cref{sec:atomic} propagates the exact constant through the HPS finite--$N$ argument and proves \Cref{thm:finite-Z}.  \Cref{sec:outlook} discusses the significance and limitations of the result and the information required for improvements beyond the cubic mean--field obstruction.  The appendix isolates the technical points that are most sensitive in a formal proof audit.
